\chapter{Introduction to eSim}

\section{Overview}

eSim (previously known as Oscad/FreeEDA) is a free and open-source Electronic Design Automation (EDA) tool developed by the FOSSEE (Free/Libre and Open Source Software for Education) project at IIT Bombay. It provides an integrated environment for circuit design, simulation, analysis, and PCB design using open-source tools such as KiCad, Ngspice, GHDL, and Verilator through a unified Python-based graphical user interface built using Python and PyQt5.

eSim serves as an accessible alternative to proprietary EDA software such as OrCAD, PSpice, and LTspice. Its open-source architecture allows students, researchers, and developers to study the internal implementation of the tool, contribute improvements, and extend functionalities according to project requirements.

The platform is supported under the National Mission on Education through ICT (NMEICT), Ministry of Education, Government of India, with the objective of promoting open-source software usage in engineering education and research.

\section{Open-Source Circuit Simulation}

The primary simulation engine used in eSim is Ngspice, which supports DC, AC, transient, and noise analysis. KiCad is integrated for schematic capture, while GHDL and Verilator enable mixed-signal simulation support through the NGHDL subsystem.

The modular and open-source nature of eSim makes it highly suitable for educational projects, internship contributions, and research-oriented development. During this internship, the same architecture was extended to integrate OpenROAD for RTL-to-GDSII backend physical design flow support.

\section{eSim Workspace and Project Structure}

The eSim workspace organizes individual circuits as separate project directories. Each project typically contains simulation files, generated netlists, schematic files, and backend outputs required during circuit analysis and implementation.

Common project files include:

\begin{itemize}
    \item \texttt{.cir} --- KiCad-generated circuit netlist
    \item \texttt{.cir.out} --- processed Ngspice-compatible netlist
    \item \texttt{.sch} --- schematic design file
    \item \texttt{.sub} --- subcircuit definitions
\end{itemize}

Figure~\ref{fig:files} shows the eSim workspace and file manager interface used during the project implementation.

\begin{figure}[H]
\centering
\includegraphics[width=0.82\textwidth]{Files.png}
\caption{eSim workspace file manager showing project structure.}
\label{fig:files}
\end{figure}

\section{Project Organisation}

The OpenROAD integration workflow was tested using example circuit projects such as \textbf{FullAdder} and \textbf{HalfAdder}. Each project directory contained generated netlists, synthesized Verilog files, backend physical design outputs, and generated layout files.

Figure~\ref{fig:faproject} shows the FullAdder project structure inside the eSim project panel.

\begin{figure}[H]
\centering
\includegraphics[width=0.50\textwidth]{FA-Project.png}
\caption{FullAdder project structure inside the eSim project panel.}
\label{fig:faproject}
\end{figure}

\section{Internship Objective}

The primary objective of this internship was to integrate the \textbf{OpenROAD} digital physical design flow into eSim, enabling users to execute a complete RTL-to-GDSII workflow directly from the eSim graphical interface.

The major objectives of the project included:

\begin{itemize}
    \item Understanding the eSim frontend and project workflow.
    \item Adding an OpenROAD execution button in the eSim graphical interface.
    \item Developing a backend Python module to execute OpenROAD Flow Scripts (ORFS).
    \item Implementing a netlist-to-RTL converter for automatic Verilog generation.
    \item Generating backend physical design outputs such as DEF and GDSII files.
    \item Testing the complete workflow using FullAdder and HalfAdder circuits.
    \item Visualizing generated layouts using OpenROAD GUI and KLayout.
\end{itemize}

\section{Project Directory Structure}

The integrated OpenROAD-eSim system follows the directory structure shown below:

\begin{lstlisting}[language=bash, caption={Project directory structure used during integration.}]
eSim/
|-- src/
|   |-- frontEnd/
|   |   `-- Application.py
|   |
|   `-- maker/
|       |-- OpenROAD.py
|       `-- netlist2rtl.py
|
|-- Examples/
|   |-- FullAdder/
|   |   |-- FullAdder.cir.out
|   |   |-- FullAdder.v
|   |   |-- FullAdder.def
|   |   |-- FullAdder.gds
|   |   `-- logs/
|   |
|   `-- HalfAdder/
|       |-- HalfAdder.cir.out
|       |-- HalfAdder.v
|       |-- HalfAdder.def
|       |-- HalfAdder.gds
|       `-- logs/
|
`-- OpenROAD-flow-scripts/
    `-- flow/
        |-- designs/
        |   `-- sky130hd/
        |       |-- FullAdder/
        |       `-- HalfAdder/
        |
        `-- results/
            `-- sky130hd/
                |-- FullAdder/
                `-- HalfAdder/
\end{lstlisting}

The above structure was used throughout the internship for OpenROAD integration, backend automation, testing, and generated file management.